\section{Discussion}
\label{sec:discussion}

The experimental results establish that the SAGE framework achieves reliable evaluation performance while revealing systematic patterns in how different quality dimensions distinguish literary categories. We organize our interpretation around the four research questions: differential discrimination patterns address discriminative capacity (RQ2), dimensional independence validates the multi-layer decomposition (RQ4), mode invariance bears on evaluation robustness (RQ3), and the multi-round architecture contributes to reliability (RQ1). We then consider broader implications for LLM evaluation methodology and AI-generated literature.

\subsection{Differential Discrimination across Interpretive Layers}

The most striking finding concerns the differential discrimination magnitudes across the three analytical layers. Cultural representation and existential-philosophical representation exhibit very large effect sizes when comparing canonical literature to LLM-generated stories ($d$=2.68 and $d$=2.40 respectively, with approximately 35\% relative gaps), while emotional-psychological representation shows a substantially smaller though still large effect ($d$=1.68, 19.1\% gap). This pattern is not an artifact of scoring range or variance: the three layers assess the same 100 stories using the same evaluation architecture, differing only in the theoretical frameworks guiding assessment. The differential must therefore reflect genuine differences in how well current generative models reproduce different dimensions of literary quality.

A plausible interpretation draws on the distinction between pattern-reproducible and stance-requiring literary capacities. Emotional structures, including affective complexity, psychological interiority, and emotional granularity, are extensively represented in narrative training corpora. Characters experience recognizable emotions; narrators convey interiority through established literary techniques; emotional arcs follow identifiable patterns. An LLM trained on large quantities of fiction can learn to reproduce these surface patterns with moderate fidelity, achieving scores around 3.36 compared to the canonical mean of 4.15. Cultural power critique and existential-philosophical inquiry, by contrast, require capacities that pattern matching alone may not deliver: a critical stance toward the social structures being represented, original philosophical engagement with questions of human finitude and meaning, and the ability to position narrative voice within contested ideological terrain. The approximately doubled effect sizes for these dimensions suggest they represent qualitatively different challenges for generative models rather than merely harder versions of the same task.

The three layers also exhibit distinct overall score distributions: emotional representation achieves the highest mean (3.96) across all text categories, followed by cultural representation (3.64) and existential-philosophical representation (3.59). This ranking suggests that emotional depth represents a more universally present quality dimension across diverse narratives, while cultural critique and philosophical engagement constitute more specialized achievements that vary substantially by genre and authorship context. The relatively high emotional scores even among LLM-generated stories support the interpretation that affective patterns are among the most learnable aspects of literary craft from training data alone.

\subsection{Canonical Literature and Genre Fiction}

The relationship between canonical literature and pulp fiction provides a complementary perspective on what the framework measures. Across all three layers, canonical works outperform pulp fiction by only modest margins (overall +0.17 points, 4.4\%), with the smallest gap appearing in emotional representation (2.7\%) and the largest in existential-philosophical representation (6.8\%). This convergence suggests that effective emotional representation constitutes a shared competence of human authorship, achieved through authentic lived experience regardless of critical positioning or institutional recognition. The larger gap in philosophical depth indicates that sustained existential inquiry distinguishes canonical literature from commercial genre fiction more strongly than either emotional or cultural sophistication, a finding consistent with the literary critical tradition that prizes philosophical ambiguity and interpretive openness as hallmarks of canonical status.

This pattern also carries methodological significance. A coarser evaluation framework that collapsed cultural, emotional, and philosophical dimensions into a single quality score would yield a modest 4.4\% canonical-pulp difference that might appear negligible. The layered decomposition reveals that this small aggregate difference masks meaningful variation: canonical and pulp fiction converge on emotional competence while diverging on philosophical depth. Such differential sensitivity across dimensions validates the ontological decomposition motivating the SAGE architecture and suggests that aggregate quality scores obscure more than they reveal.

\subsection{Dimensional Independence and Framework Validity}

The cross-layer correlation analysis (Table~\ref{tab:cross_layer_correlation}) provides direct evidence for RQ4. All three layer pairs yield moderate Pearson correlations ($r$ = 0.649--0.683), confirming that cultural, emotional-psychological, and existential-philosophical representation capture empirically distinguishable quality dimensions. The correlations are positive, as expected: higher-quality works tend to excel across multiple dimensions, and texts with impoverished cultural representation are unlikely to demonstrate profound philosophical engagement. Yet the correlations remain well below the threshold ($r>0.7$) that would suggest the layers are measuring the same underlying construct under different labels.

The asymmetry in the correlation structure is itself informative. The strongest association links cultural and existential dimensions ($r$=0.683), consistent with the observation that sustained cultural critique and philosophical inquiry both require a critical stance toward the material being represented, a capacity that distinguishes canonical literature from both commercial genre fiction and LLM-generated narratives. The weakest association links cultural and emotional dimensions ($r$=0.649), supporting the finding that emotional competence is more broadly distributed across text categories and less dependent on the analytical sophistication required for cultural power critique. These patterns reinforce the methodological claim that collapsing distinct evaluative dimensions into a single quality metric sacrifices meaningful discriminative information.

\subsection{Mode Invariance and Evaluation Grounding}

The near-perfect mode invariance observed across all layers (maximum absolute difference of 0.05 between content-limit and title-limit modes) warrants careful interpretation. In content-limit mode, the evaluator works solely from textual evidence without knowledge of authorship or critical reception. In title-limit mode, the evaluator synthesizes scholarly consensus and critical reputation without access to the text itself. That these two radically different information sources yield statistically indistinguishable assessments admits at least two readings. The optimistic interpretation holds that LLM evaluation captures genuine literary properties accessible through either direct textual analysis or accumulated critical judgment, and that the convergence validates both the framework and the consistency of the evaluator's critical reasoning. A more cautious reading notes that even content-limit evaluation cannot fully exclude the influence of training corpus knowledge, since stylistic features of well-known works may activate learned associations without explicit authorship recognition.

The truth likely involves elements of both: LLM evaluators perform genuine textual analysis while also drawing on internalized critical knowledge, producing assessments that reflect the same underlying quality dimensions regardless of the information pathway. For practical purposes, this mode invariance is advantageous. It suggests the framework can be applied with equal confidence to texts with and without established critical reception, extending its utility beyond canonical works to contemporary and emerging literature where scholarly consensus has not yet formed.

\subsection{Implications for LLM-as-Judge Methodology}

These findings contribute to the growing literature on LLM-as-judge evaluation. Prior work has demonstrated that LLM evaluators can achieve substantial agreement with human preferences in conversational and instructional settings, while also exhibiting systematic biases including position bias, verbosity bias, and self-enhancement~\cite{zheng2023judging,li2024llms}. Our results extend this literature to literary evaluation, demonstrating that with appropriate architectural design, LLM evaluators can achieve inter-rater agreement exceeding 94\% and convergence rates of 98.8\% on interpretive tasks considerably more complex than typical NLG evaluation. The multi-round architecture proves particularly important: the structured progression from initial assessment through projection bias checking, layer boundary verification, evidence sufficiency review, and final consolidation provides explicit mechanisms for addressing known sources of LLM evaluation error that single-pass approaches cannot match.

Whether this level of reliability generalizes across different LLM families and evaluation domains remains an open empirical question. The present study employs a single model (GPT-5-mini) for all evaluations. Cross-model comparison across architecturally distinct LLM families would distinguish framework-level reliability from model-specific behavioral properties and strengthen the case for multi-round iterative evaluation as a general methodology rather than a model-dependent technique.

\subsection{Implications for AI-Generated Literature}

The systematic capability profile revealed by the framework carries implications beyond evaluation methodology. The finding that current LLMs achieve moderate competence at emotional pattern reproduction while demonstrating fundamental weaknesses in cultural power critique and existential-philosophical depth provides empirically grounded guidance for both AI development and literary assessment. For AI developers, these results indicate that improvements in surface fluency and emotional mimicry should not be mistaken for progress toward genuine literary capability, since the most distinctive achievements of canonical literature lie precisely in the dimensions where AI generation performs worst. For literary scholars and critics engaging with AI-generated content, the framework provides systematic criteria for distinguishing technically competent but philosophically shallow AI narratives from works exhibiting the critical depth and existential authenticity associated with the human literary tradition.

More broadly, the approximately 35\% relative gaps in cultural and existential dimensions raise the question of whether these weaknesses reflect fundamental architectural limitations of current language models or temporary gaps addressable through improved training and prompting strategies. The pattern-reproducibility interpretation developed above suggests the former: if cultural critique and philosophical depth require genuine critical stance and existential confrontation rather than statistical pattern matching, then scaling training data or refining prompts may yield diminishing returns. Longitudinal evaluation as generative models improve would provide empirical evidence to adjudicate this question, tracking whether the capability hierarchy identified here proves stable or shifts with architectural advances.
